% =====================================================================
%  chapters/minggu-16.tex
%  Minggu 16 — Presentasi, Refleksi & Finalisasi Proyek Akhir
% =====================================================================

\chapter{Presentasi, Refleksi \& Finalisasi Proyek Akhir}
\label{chap:minggu-16}

% ---------------------------------------------------------------------
% 1. CAPAIAN PEMBELAJARAN
% ---------------------------------------------------------------------
\begin{capaian}
Setelah menyelesaikan bab ini, mahasiswa diharapkan mampu:
\begin{itemize}
  \item mempresentasikan proses proyek redesign dari riset $\rightarrow$ desain $\rightarrow$ implementasi $\rightarrow$ deployment;
  \item menjawab pertanyaan tentang kode dan keputusan desain secara individual;
  \item merefleksikan perkembangan diri selama semester dalam mata kuliah ini;
  \item menyelesaikan seluruh \emph{deliverable} proyek akhir.
\end{itemize}
\end{capaian}

% ---------------------------------------------------------------------
% 2. TUJUAN PRAKTIK
% ---------------------------------------------------------------------
\begin{tujuanpraktik}
Pada sesi praktik ini, mahasiswa akan:
\begin{itemize}
  \item melakukan pengecekan kesiapan teknis sebelum presentasi;
  \item mempresentasikan hasil proyek redesign di depan kelas;
  \item menjawab pertanyaan individual dari pengajar;
  \item mengisi refleksi diri akhir semester;
  \item melakukan finalisasi \emph{repository}.
\end{itemize}
\end{tujuanpraktik}

% ---------------------------------------------------------------------
% 3. PERALATAN
% ---------------------------------------------------------------------
\section{Peralatan yang Dibutuhkan}
\label{sec:peralatan-16}

\begin{itemize}
  \item \textbf{Browser} --- \emph{live demo} website.
  \item \textbf{GitHub} --- menunjukkan \emph{repository} dan \emph{commit history}.
  \item \textbf{Figma} --- menunjukkan file desain.
  \item \textbf{Slide presentasi} --- pembukaan dan penutup (Google Slides / Markdown).
  \item \textbf{VS Code} --- menunjukkan kode saat sesi tanya-jawab.
\end{itemize}

% ---------------------------------------------------------------------
% 4. PERSIAPAN SEBELUM PRESENTASI
% ---------------------------------------------------------------------
\section{Persiapan Sebelum Presentasi}
\label{sec:persiapan-sebelum-presentasi}

\subsection{Cek Kesiapan Teknis}
\label{subsec:cek-kesiapan}

\begin{langkah}
  \item Buka website via GitHub Pages di browser --- coba navigasi seluruh halaman.
  \item \emph{Resize} browser dari desktop ke mobile --- pastikan tampilan responsif.
  \item Pastikan \emph{repository} GitHub bersih --- tidak ada file \emph{temporary} atau file yang tidak perlu.
  \item Pastikan \texttt{README.md} sudah lengkap.
  \item Pastikan semua anggota bisa mengakses \emph{repository} dari laptop masing-masing.
\end{langkah}

\subsection{Siapkan Tab Browser}
\label{subsec:siapkan-tab}

Buka semua tab berikut di browser sebelum giliran presentasi:
\begin{enumerate}
  \item Tab 1: Website hasil redesign (GitHub Pages).
  \item Tab 2: Repository GitHub (terlihat \emph{commit history} \& README).
  \item Tab 3: File desain Figma.
  \item Tab 4: File \emph{heuristic evaluation} di GitHub (atau VS Code).
\end{enumerate}

\begin{tangkapanlayar}
  {assets/images/minggu-16/persiapan-tab-browser.png}
  {Persiapan presentasi --- browser dengan 4 tab terbuka.}
  {Buka browser dan pastikan keempat tab sudah siap sebelum giliran presentasi kelompok Anda.}
\end{tangkapanlayar}

% ---------------------------------------------------------------------
% 5. PRESENTASI KELOMPOK
% ---------------------------------------------------------------------
\section{Presentasi Kelompok}
\label{sec:presentasi}

Gunakan \emph{outline} yang sudah disiapkan di minggu 15. Waktu per kelompok: 8--10 menit presentasi + 5 menit tanya-jawab.

\subsection{Pembukaan (± 1 Menit)}
\label{subsec:pembukaan}

\begin{itemize}
  \item Perkenalan anggota kelompok.
  \item Nama website yang diriset dan URL-nya.
  \item URL website redesign (GitHub Pages).
\end{itemize}

\begin{tangkapanlayar}
  {assets/images/minggu-16/slide-pembukaan.png}
  {Slide pembukaan --- nama kelompok, nama website.}
  {Siapkan slide pembukaan yang memuat nama kelompok dan nama website yang diredesign.}
\end{tangkapanlayar}

\subsection{Riset UI/UX (± 2--3 Menit)}
\label{subsec:presentasi-riset}

Tampilkan:
\begin{enumerate}
  \item \emph{Screenshot} halaman asli website yang diriset.
  \item Top 3 \emph{pain point} dari \emph{heuristic evaluation} --- jelaskan apa masalahnya dan mengapa itu masalah bagi pengguna.
  \item 1--2 temuan \emph{benchmarking} --- apa yang dipelajari dari website lain.
\end{enumerate}

\begin{catatan}
Contoh penyampaian: ``Di halaman asli, tombol `Ajukan Permohonan' tidak terlihat karena warnanya menyatu dengan \emph{background}. Ini melanggar heuristic \#1 --- Visibility of System Status.''
\end{catatan}

\begin{tangkapanlayar}
  {assets/images/minggu-16/slide-riset-pain-point.png}
  {Tampilan slide riset --- pain point utama.}
  {Pastikan slide riset memuat \emph{screenshoot} halaman asli dan daftar \emph{pain point} utama.}
\end{tangkapanlayar}

\subsection{Desain Ulang (± 2--3 Menit)}
\label{subsec:presentasi-desain}

Tampilkan:
\begin{enumerate}
  \item \emph{Mockup} Figma (desktop \& mobile) --- tunjukkan langsung di Figma.
  \item 3 keputusan desain utama, masing-masing dijelaskan dengan kaitannya terhadap temuan riset.
\end{enumerate}

\begin{tip}
Contoh penyampaian: ``Keputusan 1: Kami memindahkan tombol CTA ke posisi lebih prominent. Ini karena di riset ditemukan bahwa tombol utama sulit ditemukan (\emph{pain point} \#1).''
\end{tip}

\begin{tangkapanlayar}
  {assets/images/minggu-16/mockup-figma-desktop.png}
  {Mockup Figma --- tampilan desktop.}
  {Tunjukkan \emph{mockup} desktop di Figma, jelaskan keputusan desain utama.}
\end{tangkapanlayar}

\begin{tangkapanlayar}
  {assets/images/minggu-16/mockup-figma-mobile.png}
  {Mockup Figma --- tampilan mobile.}
  {Tunjukkan \emph{mockup} mobile di Figma, jelaskan bagaimana desain beradaptasi.}
\end{tangkapanlayar}

\subsection{Live Demo (± 2 Menit)}
\label{subsec:live-demo}

\begin{langkah}
  \item Buka website di browser (GitHub Pages).
  \item Tunjukkan fitur utama: navigasi, konten utama, responsif (resize browser dari desktop ke mobile).
  \item Tunjukkan \emph{commit history} di GitHub --- siapa mengerjakan apa.
\end{langkah}

\begin{tangkapanlayar}
  {assets/images/minggu-16/live-demo-desktop.png}
  {Live demo website --- tampilan desktop.}
  {Buka website di browser, tunjukkan navigasi dan konten utama.}
\end{tangkapanlayar}

\begin{tangkapanlayar}
  {assets/images/minggu-16/live-demo-mobile.png}
  {Live demo website --- tampilan mobile.}
  {Resize browser ke ukuran mobile, tunjukkan responsivitas.}
\end{tangkapanlayar}

\begin{tangkapanlayar}
  {assets/images/minggu-16/commit-history-github.png}
  {Commit history di GitHub.}
  {Buka tab GitHub, tunjukkan \emph{commit history} untuk membuktikan kontribusi per anggota.}
\end{tangkapanlayar}

\subsection{Pembagian Kerja dan Penggunaan AI (± 1 Menit)}
\label{subsec:pembagian-kerja-presentasi}

\begin{itemize}
  \item Tampilkan tabel pembagian tugas beserta jumlah \emph{commit} per anggota.
  \item Jelaskan tools apa yang dipakai (jika ada), di bagian mana, dan tunjukkan dokumentasi \emph{prompt}-nya.
\end{itemize}

\subsection{Penutup (± 1 Menit)}
\label{subsec:penutup-presentasi}

\begin{itemize}
  \item Satu hal terpenting yang dipelajari dari proyek ini.
  \item Satu tantangan terbesar yang dihadapi.
\end{itemize}

% ---------------------------------------------------------------------
% 6. TANYA-JAWAB INDIVIDUAL
% ---------------------------------------------------------------------
\section{Tanya-Jawab Individual}
\label{sec:tanya-jawab}

Setelah presentasi, pengajar mengajukan pertanyaan kepada setiap anggota secara individual ($\pm$ 5 menit per kelompok). Berikut contoh pertanyaan yang mungkin ditanyakan berdasarkan bagian yang dikerjakan:

\begin{table}[htbp]
\centering
\caption{Contoh Pertanyaan Tanya-Jawab Individual}
\label{tab:pertanyaan-tj}
\small
\begin{tabular}{@{}p{5cm}p{5cm}p{5cm}@{}}
\toprule
\textbf{Header/Navigasi} & \textbf{Konten Utama} & \textbf{Footer/Responsive} \\
\midrule
Bagaimana cara kerja navigasi di website ini? Jelaskan HTML-nya. & Kenapa formulir ini diletakkan di posisi ini? & Bagaimana cara \emph{media query} bekerja di CSS ini? \\
Apa heuristik yang paling relevan dengan perubahan di header? & Bagaimana website ini mengatasi \emph{pain point} \#2? & Apa bedanya \texttt{flex-direction: row} dan \texttt{column} di kode ini? \\
Jika saya ingin mengubah warna header, baris CSS mana yang perlu diubah? & Bisakah kamu menjelaskan alur pengguna mengisi formulir ini? & Di mana batas \emph{breakpoint} untuk mobile di CSS ini? \\
\bottomrule
\end{tabular}
\end{table}

\begin{tip}
Untuk mahasiswa: Buka VS Code saat sesi tanya-jawab --- tunjukkan kode langsung. Jujur jika tidak tahu, tapi tunjukkan usaha untuk mencari tahu. Hubungkan jawaban dengan riset --- misal: ``Ini kami buat karena di riset ditemukan\ldots''
\end{tip}

\begin{tangkapanlayar}
  {assets/images/minggu-16/tanya-jawab-vscode.png}
  {Sesi tanya-jawab --- anggota menjelaskan kode di VS Code.}
  {Buka VS Code di layar, navigasi ke file yang relevan dengan pertanyaan pengajar.}
\end{tangkapanlayar}

% ---------------------------------------------------------------------
% 7. REFLEKSI INDIVIDU
% ---------------------------------------------------------------------
\section{Refleksi Individu}
\label{sec:refleksi}

Setelah semua kelompok selesai presentasi, setiap mahasiswa mengisi refleksi diri. Buka file \texttt{reflection.md} di \emph{repo} (atau gunakan form \emph{online} yang disediakan pengajar) dan jawab pertanyaan-pertanyaan berikut.

\subsection{Perkembangan Skill}
\label{subsec:perkembangan-skill}

Buat tabel perbandingan perkembangan skill dari awal semester ke akhir semester untuk kompetensi: HTML, CSS, JavaScript, Figma, Git/GitHub, dan \emph{Heuristic Evaluation}.

\subsection{Refleksi Proyek Akhir}
\label{subsec:refleksi-proyek}

Jawab pertanyaan-pertanyaan berikut secara jujur dan reflektif:
\begin{enumerate}
  \item Apa yang paling menantang dari proyek ini?
  \item Apa yang paling memuaskan dari hasil kerja kelompok?
  \item Jika bisa mengulang, apa yang akan dilakukan berbeda?
  \item Bagian mana dari kode yang paling dipahami dan bisa dijelaskan dengan baik?
  \item Bagian mana yang masih perlu dipelajari lebih lanjut?
\end{enumerate}

\subsection{Refleksi Pembelajaran}
\label{subsec:refleksi-pembelajaran}

Jawab pertanyaan-pertanyaan berikut:
\begin{enumerate}
  \item Satu hal terpenting yang dipelajari di mata kuliah ini.
  \item Bagaimana mata kuliah ini mengubah cara pandang terhadap \emph{web development}.
  \item Apakah ada topik yang ingin dipelajari lebih lanjut?
\end{enumerate}

\begin{tangkapanlayar}
  {assets/images/minggu-16/refleksi-terisi.png}
  {Formulir refleksi yang sudah diisi.}
  {Buka \texttt{reflection.md} di VS Code, pastikan seluruh pertanyaan sudah dijawab.}
\end{tangkapanlayar}

\begin{langkah}
  \item Isi seluruh pertanyaan refleksi di \texttt{reflection.md}.
  \item \emph{Commit} dan \emph{push}:
\begin{lstlisting}[language=bash, caption={Commit refleksi}]
git add reflection.md
git commit -m "docs: refleksi akhir semester - [nama]"
git push
\end{lstlisting}
\end{langkah}

% ---------------------------------------------------------------------
% 8. FINALISASI REPOSITORY
% ---------------------------------------------------------------------
\section{Finalisasi Repository}
\label{sec:finalisasi-repo}

Pastikan \emph{repository} dalam kondisi final dan siap untuk portofolio.

\subsection{Review Checklist Final}
\label{subsec:checklist-final}

\begin{langkah}
  \item Pastikan \texttt{index.html} berfungsi dengan benar.
  \item Pastikan CSS responsif (desktop \& mobile).
  \item Pastikan JavaScript berfungsi (jika ada).
  \item Pastikan gambar \& aset semua tampil.
  \item Pastikan \texttt{README.md} lengkap.
  \item Pastikan folder \texttt{research/} berisi \emph{heuristic evaluation} \& \emph{benchmarking}.
  \item Pastikan file Figma sudah di-\emph{share} (tautan di README).
  \item Pastikan website ter-deploy di GitHub Pages.
  \item Pastikan \emph{commit history} mencerminkan kontribusi per anggota.
  \item Pastikan tidak ada file \emph{temporary} (\texttt{.DS\_Store}, \texttt{Thumbs.db}, dsb).
\end{langkah}

\subsection{Tambahkan File .gitignore}
\label{subsec:gitignore}

Buat \texttt{.gitignore} di \emph{root repository} jika belum ada:

\begin{lstlisting}[caption={Contoh file .gitignore}]
# OS files
.DS_Store
Thumbs.db

# Editor files
.vscode/settings.json
*.swp
*.swo

# Node modules (jika menggunakan Tailwind)
node_modules/

# Build output
dist/
\end{lstlisting}

\subsection{Final Commit}
\label{subsec:final-commit}

\begin{lstlisting}[language=bash, caption={Final commit proyek akhir}]
git add .
git commit -m "chore: finalisasi repository proyek akhir"
git push
\end{lstlisting}

\begin{tangkapanlayar}
  {assets/images/minggu-16/repository-final-github.png}
  {Repository GitHub --- tampilan root yang bersih dan lengkap.}
  {Buka \emph{repository} di GitHub, pastikan seluruh struktur sudah sesuai dan tidak ada file yang tidak perlu.}
\end{tangkapanlayar}

\begin{tangkapanlayar}
  {assets/images/minggu-16/commit-history-akhir.png}
  {Commit history akhir --- terlihat kontribusi merata per anggota.}
  {Buka tab Insights $\rightarrow$ Contributors, pastikan kontribusi merata di antara anggota.}
\end{tangkapanlayar}

% ---------------------------------------------------------------------
% 9. STRUKTUR REPOSITORY FINAL
% ---------------------------------------------------------------------
\section{Struktur Repository Final}
\label{sec:struktur-final}

Berikut struktur \emph{repository} yang diharapkan setelah proyek akhir selesai:

\begin{lstlisting}[caption={Struktur folder repository final}]
nama-repo/
  README.md
  .gitignore
  research/
    heuristic-evaluation.md
    benchmarking.md
    rasional-desain.md
  src/
    index.html
    css/
      style.css
    js/
      script.js
    assets/
      logo.png
  design/
    figma-link.txt
  docs/
    presentasi-outline.md
  reflection.md
\end{lstlisting}

% ---------------------------------------------------------------------
% 10. DELIVERABLE AKHIR PROYEK
% ---------------------------------------------------------------------
\section{Deliverable Akhir Proyek}
\label{sec:deliverable}

Pastikan seluruh \emph{deliverable} berikut sudah selesai:
\begin{itemize}
  \item Dokumen hasil riset UI/UX (\texttt{research/heuristic-evaluation.md}).
  \item Dokumen \emph{benchmarking} (\texttt{research/benchmarking.md}).
  \item File Figma (\emph{wireframe} + \emph{mockup}) --- tautan di README.
  \item Repository GitHub (\emph{source code} + README).
  \item Website ter-deploy (GitHub Pages) --- tautan di README.
  \item Pembagian tugas di README (\texttt{README.md}).
  \item Dokumentasi penggunaan \emph{agentic coding} (jika ada).
  \item Refleksi individu (\texttt{reflection.md}).
\end{itemize}

% ---------------------------------------------------------------------
% 11. TROUBLESHOOTING
% ---------------------------------------------------------------------
\section{Troubleshooting}
\label{sec:troubleshooting-16}

\begin{table}[htbp]
\centering
\caption{Masalah Umum dan Solusi --- Minggu 16}
\label{tab:troubleshooting-16}
\small
\begin{tabular}{@{}p{5.5cm}p{8cm}@{}}
\toprule
\textbf{Masalah} & \textbf{Solusi} \\
\midrule
Website tidak bisa diakses saat presentasi & Buka file \texttt{index.html} langsung dari VS Code (Live Server) sebagai cadangan. Atau unduh \emph{repo} dan buka secara lokal. \\
Figma tidak bisa dibuka & Pastikan akun Figma sudah \emph{login}. Jika di lab kampus, minta bantuan admin untuk install Figma \emph{desktop}. \\
Lupa \emph{commit} bagian tertentu & Cek \texttt{git status} $\rightarrow$ jika ada file yang belum ter-\emph{commit}, lakukan \emph{commit} sekarang. Jangan lupa \emph{push}. \\
Tidak sempat mengisi refleksi & Refleksi bisa diisi dan di-\emph{push} setelah kelas. Deadline: $1\times24$ jam setelah presentasi. \\
Koneksi internet mati saat presentasi & Gunakan versi lokal (Live Server) sebagai cadangan. Pastikan sudah \emph{download} \emph{repo} sebelumnya. \\
Anggota tidak bisa hadir & Hubungi pengajar sebelum kelas. Kemungkinan presentasi dijadwalkan ulang atau dilakukan \emph{remote}. \\
Website \emph{deploy} expired/blank & Cek Settings $\rightarrow$ Pages di GitHub. Pastikan \emph{branch source} masih aktif. \emph{Push} ulang jika perlu. \\
\emph{Merge conflict} di final \emph{commit} & Selesaikan conflict: buka file $\rightarrow$ pilih versi yang benar $\rightarrow$ hapus marker \texttt{<<<<<<<} $\rightarrow$ \emph{commit}. \\
\bottomrule
\end{tabular}
\end{table}

% ---------------------------------------------------------------------
% 12. TAUTAN REFERENSI
% ---------------------------------------------------------------------
\section{Tautan Referensi}
\label{sec:referensi-minggu-16}

\begin{itemize}
  \item GitHub Pages: \url{https://docs.github.com/en/pages}
  \item Nielsen Norman Group: \url{https://www.nngroup.com/articles/ten-usability-heuristics/}
  \item Figma: \url{https://www.figma.com}
  \item MDN Web Docs: \url{https://developer.mozilla.org}
  \item GitHub Skills: \url{https://skills.github.com}
\end{itemize}

% ---------------------------------------------------------------------
% 13. LATIHAN
% ---------------------------------------------------------------------
\section{Latihan Mandiri}
\label{sec:latihan-minggu-16}

\begin{latihan}[Refleksi dan Portofolio]
\begin{enumerate}[label=\alph*., leftmargin=2em, itemsep=2pt]
  \item Buka \emph{repository} proyek akhir Anda di GitHub. Tinjau \emph{commit history} dan identifikasi minggu mana yang memiliki kontribusi terbanyak. Apakah ada pola yang bisa diperbaiki di proyek berikutnya?
  \item Buka website hasil redesign Anda di dua browser berbeda (misal Chrome dan Firefox). Catat minimal 1 perbedaan tampilan dan jelaskan penyebabnya.
\end{enumerate}
\end{latihan}
